% TODO checklist: Inspected files for this section:
% - Q&A.md (questions 46, 50 - recommendations, roadmap priorities)
% - comparison.md (integration patterns, what leaders offer)
% - README.md (roadmap hints, architecture extensibility)
% - sapthesis-doc.tex (future work topics list)

\section{Future Work}
\label{sec:future-work}

The Cineca Agentic Platform provides a solid foundation for enterprise AI agent orchestration, but significant opportunities exist for enhancement and extension. This section outlines prioritized directions for future development, organized by impact and feasibility.

\subsection{High-Priority Enhancements}

\paragraph{1. Real-Time Streaming Responses}

The current polling-based architecture for agent run updates introduces latency and increases backend load. Implementing true push-based streaming would significantly improve user experience:

\begin{itemize}
    \item \textbf{WebSocket integration}: Persistent connections for bidirectional communication
    \item \textbf{Step-by-step streaming}: Real-time emission of orchestration steps as they complete
    \item \textbf{Token streaming}: Character-by-character LLM output for responsive chat interfaces
    \item \textbf{Progress indicators}: Fine-grained progress updates for long-running operations
\end{itemize}

Implementation considerations:
\begin{itemize}
    \item WebSocket connection management at scale (connection pooling, load balancing)
    \item Fallback to SSE/polling for clients that don't support WebSockets
    \item Authentication and tenant isolation for persistent connections
    \item Integration with existing observability infrastructure
\end{itemize}

\paragraph{2. Conversation Memory Systems}

Extending beyond session-based state to implement sophisticated memory systems would enable more capable long-running agent interactions:

\begin{itemize}
    \item \textbf{Summary memory}: Automatic compression of conversation history to fit context windows
    \item \textbf{Semantic memory}: Embedding-based storage and retrieval of important conversation points
    \item \textbf{Entity memory}: Tracking and recall of entities mentioned across sessions
    \item \textbf{Episodic memory}: Structured recall of past interactions for continuity
\end{itemize}

This enhancement would leverage the existing Memgraph infrastructure, storing memory as graph nodes and relationships with temporal metadata.

\paragraph{3. Document Ingestion and RAG Pipeline}

Complementing the graph-native NL-to-Cypher pipeline with document-based retrieval would significantly expand the platform's applicability:

\begin{itemize}
    \item \textbf{Document loaders}: Support for PDF, Word, HTML, Markdown, and structured data formats
    \item \textbf{Chunking strategies}: Semantic, sentence-based, and recursive chunking with overlap
    \item \textbf{Embedding generation}: Integration with embedding providers (OpenAI, Cohere, local models)
    \item \textbf{Vector storage}: PostgreSQL pgvector extension or dedicated vector database integration
    \item \textbf{Hybrid retrieval}: Combining vector similarity with graph traversal for comprehensive answers
\end{itemize}

The multi-tenant architecture would naturally extend to vector storage, with per-tenant embedding namespaces.

\subsection{Agent Capability Extensions}

\paragraph{4. Multi-Agent Collaboration Framework}

Enabling multiple agents to collaborate on complex tasks would significantly expand the platform's capabilities:

\begin{itemize}
    \item \textbf{Agent-to-agent communication}: Message passing between specialized agents
    \item \textbf{Task delegation}: Hierarchical agent structures with manager and worker agents
    \item \textbf{Shared context}: Common knowledge base accessible to all collaborating agents
    \item \textbf{Coordination protocols}: Turn-taking, voting, and consensus mechanisms
\end{itemize}

Implementation approaches:
\begin{itemize}
    \item Extend the orchestrator to manage agent teams rather than single agents
    \item Define agent role specifications (planner, researcher, executor, critic)
    \item Implement conversation routing between agents
    \item Maintain audit trails showing agent interactions
\end{itemize}

\paragraph{5. Additional Agent Reasoning Patterns}

Beyond the current ReAct-style orchestration, implementing alternative reasoning patterns would enable more sophisticated problem-solving:

\begin{itemize}
    \item \textbf{Plan-and-Execute}: Separate planning and execution phases with plan revision
    \item \textbf{Tree-of-Thought}: Exploration of multiple reasoning paths with backtracking
    \item \textbf{Self-Ask}: Decomposition into sub-questions with independent resolution
    \item \textbf{Reflexion}: Self-critique and iterative refinement of responses
    \item \textbf{ReWOO}: Reasoning without observation for reduced LLM calls
\end{itemize}

The orchestrator would become a pattern dispatcher, selecting the appropriate reasoning strategy based on task characteristics.

\paragraph{6. Enhanced NL-to-Cypher Capabilities}

Extending the NL-to-Cypher pipeline to support more complex query types:

\begin{itemize}
    \item \textbf{Multi-hop reasoning}: Complex queries requiring multiple graph traversals
    \item \textbf{Temporal queries}: Time-based filtering and trend analysis
    \item \textbf{Aggregation combinations}: Nested GROUP BY with multiple aggregation functions
    \item \textbf{Subgraph extraction}: Returning connected subgraphs rather than tabular results
    \item \textbf{Write operations}: Controlled, audited graph modifications through natural language
\end{itemize}

Safety considerations for write operations would require additional validation layers and approval workflows.

\subsection{Infrastructure and Operations}

\paragraph{7. Distributed Worker Architecture}

Scaling the job processing infrastructure for high-volume deployments:

\begin{itemize}
    \item \textbf{Distributed locking}: Redis SETNX or Redlock for multi-worker coordination
    \item \textbf{Worker pools}: Dynamic scaling of worker processes based on queue depth
    \item \textbf{Job priority queues}: Multiple priority levels with fair scheduling
    \item \textbf{Dead letter queues}: Automatic routing of repeatedly failed jobs for investigation
    \item \textbf{Job retry strategies}: Configurable exponential backoff with jitter
\end{itemize}

\paragraph{8. Event Sourcing Architecture}

Implementing event sourcing for complete audit trails and operational flexibility:

\begin{itemize}
    \item \textbf{Event store}: Append-only log of all state changes
    \item \textbf{State reconstruction}: Ability to rebuild current state from event history
    \item \textbf{Temporal queries}: Query historical states at any point in time
    \item \textbf{Event replay}: Debugging and testing by replaying production events
    \item \textbf{Integration hooks}: External system integration via event consumption
\end{itemize}

This architecture would enhance compliance capabilities and enable sophisticated debugging scenarios.

\paragraph{9. Microservices Decomposition}

For larger deployments, decomposing the monolithic backend into specialized services:

\begin{itemize}
    \item \textbf{Agent service}: Orchestration and run management
    \item \textbf{Tool service}: MCP tool registry and invocation
    \item \textbf{Auth service}: Authentication and authorization
    \item \textbf{Graph service}: Memgraph operations and NL-to-Cypher
    \item \textbf{Job service}: Background job processing
    \item \textbf{Metrics service}: Observability aggregation
\end{itemize}

This decomposition would enable:
\begin{itemize}
    \item Independent scaling of high-demand components
    \item Technology diversity (different languages/frameworks per service)
    \item Smaller deployment units for faster iteration
    \item Team ownership boundaries
\end{itemize}

\subsection{LLM Provider Enhancements}

\paragraph{10. Additional LLM Provider Integrations}

Expanding the provider ecosystem beyond Ollama and OpenAI-compatible APIs:

\begin{itemize}
    \item \textbf{Anthropic Claude}: Direct API integration with Claude models
    \item \textbf{Google Gemini}: Integration with Gemini API
    \item \textbf{AWS Bedrock}: Multi-model access through AWS
    \item \textbf{Azure OpenAI}: Enterprise Azure deployments
    \item \textbf{Cohere}: Command and Embed model families
    \item \textbf{Local inference}: vLLM, text-generation-inference for self-hosted deployment
\end{itemize}

Each provider integration would include:
\begin{itemize}
    \item Authentication configuration
    \item Rate limiting adaptation
    \item Cost tracking calibration
    \item Circuit breaker thresholds
    \item Model-specific prompt formatting
\end{itemize}

\paragraph{11. Advanced Cost Governance}

Extending the cost tracking framework for enterprise budget management:

\begin{itemize}
    \item \textbf{Budget hierarchies}: Organization $\rightarrow$ department $\rightarrow$ team $\rightarrow$ user
    \item \textbf{Soft and hard limits}: Warnings vs. enforcement thresholds
    \item \textbf{Cost allocation}: Chargeback reports for internal billing
    \item \textbf{Predictive alerts}: Forecasting budget exhaustion based on trends
    \item \textbf{Provider cost optimization}: Automatic routing to cost-effective providers
\end{itemize}

\paragraph{12. Model Warm-up and Preloading}

Reducing cold-start latency for local LLM deployments:

\begin{itemize}
    \item \textbf{Startup warm-up}: Automatic model loading on service initialization
    \item \textbf{Model preloading}: Keep frequently-used models in memory
    \item \textbf{Predictive loading}: Anticipate model needs based on usage patterns
    \item \textbf{Memory management}: Intelligent model eviction based on usage recency
\end{itemize}

\subsection{Security and Compliance}

\paragraph{13. Enhanced PII Protection}

Extending data protection capabilities:

\begin{itemize}
    \item \textbf{ML-based detection}: Named entity recognition for PII beyond regex patterns
    \item \textbf{Contextual redaction}: Preserve meaning while removing identifiers
    \item \textbf{Reversible tokenization}: Secure tokens for authorized re-identification
    \item \textbf{Cross-border compliance}: Region-aware data handling (GDPR, CCPA)
    \item \textbf{Differential privacy}: Aggregate query protection for analytics
\end{itemize}

\paragraph{14. Advanced Audit and Compliance}

Strengthening audit capabilities for regulated environments:

\begin{itemize}
    \item \textbf{Tamper-evident logs}: Cryptographic chaining of audit records
    \item \textbf{Compliance reports}: Automated generation of compliance documentation
    \item \textbf{Access certification}: Periodic review and attestation workflows
    \item \textbf{Data lineage}: Tracking of information flow through the system
    \item \textbf{Retention policies}: Automated enforcement of data retention requirements
\end{itemize}

\subsection{User Experience}

\paragraph{15. Unified UI Framework}

Addressing the dual-UI maintenance burden:

\begin{itemize}
    \item \textbf{Shared component library}: Common UI components for both interfaces
    \item \textbf{API SDK generation}: Auto-generated clients for TypeScript and Python
    \item \textbf{Design system}: Consistent visual language across interfaces
    \item \textbf{Embedded panels}: Integration of control panel features into chat UI
\end{itemize}

\paragraph{16. Agent Debugging and Visualization}

Enhanced tooling for understanding agent behavior:

\begin{itemize}
    \item \textbf{Step visualization}: Graphical representation of orchestration flow
    \item \textbf{Decision tree display}: Visualization of intent classification paths
    \item \textbf{Token usage analysis}: Per-step token consumption breakdown
    \item \textbf{Replay functionality}: Re-execute past runs with modified inputs
    \item \textbf{Comparison views}: Side-by-side analysis of different run configurations
\end{itemize}

\subsection{Research Directions}

\paragraph{17. Benchmark Development}

Creating standardized benchmarks for enterprise agent platforms:

\begin{itemize}
    \item \textbf{Security benchmarks}: Tenant isolation, injection resistance, authorization correctness
    \item \textbf{Performance benchmarks}: Latency, throughput, scalability characteristics
    \item \textbf{Reliability benchmarks}: Failure recovery, degradation behavior
    \item \textbf{Capability benchmarks}: Task completion rates across domain-specific scenarios
\end{itemize}

\paragraph{18. Federated Multi-Cluster Deployment}

Supporting distributed deployments across multiple data centers:

\begin{itemize}
    \item \textbf{Data sovereignty}: Tenant data restricted to specific regions
    \item \textbf{Cross-cluster routing}: Request routing based on tenant affinity
    \item \textbf{Replicated configuration}: Synchronized tenant and model configuration
    \item \textbf{Global observability}: Unified metrics and logging across clusters
\end{itemize}

\subsection{Integration Ecosystem}

\paragraph{19. External Workflow Engine Integration}

Connecting with enterprise workflow systems:

\begin{itemize}
    \item \textbf{Temporal integration}: Long-running workflows with Temporal
    \item \textbf{Airflow operators}: Agent runs as Airflow tasks
    \item \textbf{Prefect tasks}: Agent orchestration within Prefect flows
    \item \textbf{Kubernetes Jobs}: Native Kubernetes job integration
\end{itemize}

\paragraph{20. MCP Protocol Interoperability}

Full compliance with the Model Context Protocol specification:

\begin{itemize}
    \item \textbf{MCP server adapter}: Expose platform tools through standard MCP interface
    \item \textbf{MCP client integration}: Consume external MCP servers as tool providers
    \item \textbf{Tool discovery}: Dynamic discovery of MCP servers in the network
    \item \textbf{Schema translation}: Automatic mapping between MCP schemas and internal formats
\end{itemize}

This would enable the platform to participate in the broader MCP ecosystem, both as a provider and consumer of MCP-compliant tools.

\subsection{Roadmap Prioritization}

Based on impact and feasibility, the following prioritization is recommended:

\begin{table}[ht]
\centering
\caption{Future work prioritization matrix.}
\label{tab:roadmap-priority}
\small
\begin{tabular}{p{5cm}cc}
\hline
\textbf{Enhancement} & \textbf{Impact} & \textbf{Effort} \\
\hline
Real-time streaming & High & Medium \\
Conversation memory & High & Medium \\
Distributed workers & High & Low \\
Multi-agent collaboration & High & High \\
Document RAG pipeline & Medium & High \\
Additional LLM providers & Medium & Low \\
Event sourcing & Medium & High \\
MCP interoperability & Medium & Medium \\
\hline
\end{tabular}
\end{table}

The recommended sequence prioritizes high-impact, lower-effort enhancements (streaming, distributed workers, LLM providers) before tackling larger architectural changes (multi-agent, event sourcing, microservices decomposition).

